% Options for packages loaded elsewhere
% Options for packages loaded elsewhere
\PassOptionsToPackage{unicode}{hyperref}
\PassOptionsToPackage{hyphens}{url}
\PassOptionsToPackage{dvipsnames,svgnames,x11names}{xcolor}
%
\documentclass[
  letterpaper,
  DIV=11,
  numbers=noendperiod]{scrartcl}
\usepackage{xcolor}
\usepackage{amsmath,amssymb}
\setcounter{secnumdepth}{-\maxdimen} % remove section numbering
\usepackage{iftex}
\ifPDFTeX
  \usepackage[T1]{fontenc}
  \usepackage[utf8]{inputenc}
  \usepackage{textcomp} % provide euro and other symbols
\else % if luatex or xetex
  \usepackage{unicode-math} % this also loads fontspec
  \defaultfontfeatures{Scale=MatchLowercase}
  \defaultfontfeatures[\rmfamily]{Ligatures=TeX,Scale=1}
\fi
\usepackage{lmodern}
\ifPDFTeX\else
  % xetex/luatex font selection
\fi
% Use upquote if available, for straight quotes in verbatim environments
\IfFileExists{upquote.sty}{\usepackage{upquote}}{}
\IfFileExists{microtype.sty}{% use microtype if available
  \usepackage[]{microtype}
  \UseMicrotypeSet[protrusion]{basicmath} % disable protrusion for tt fonts
}{}
\makeatletter
\@ifundefined{KOMAClassName}{% if non-KOMA class
  \IfFileExists{parskip.sty}{%
    \usepackage{parskip}
  }{% else
    \setlength{\parindent}{0pt}
    \setlength{\parskip}{6pt plus 2pt minus 1pt}}
}{% if KOMA class
  \KOMAoptions{parskip=half}}
\makeatother
% Make \paragraph and \subparagraph free-standing
\makeatletter
\ifx\paragraph\undefined\else
  \let\oldparagraph\paragraph
  \renewcommand{\paragraph}{
    \@ifstar
      \xxxParagraphStar
      \xxxParagraphNoStar
  }
  \newcommand{\xxxParagraphStar}[1]{\oldparagraph*{#1}\mbox{}}
  \newcommand{\xxxParagraphNoStar}[1]{\oldparagraph{#1}\mbox{}}
\fi
\ifx\subparagraph\undefined\else
  \let\oldsubparagraph\subparagraph
  \renewcommand{\subparagraph}{
    \@ifstar
      \xxxSubParagraphStar
      \xxxSubParagraphNoStar
  }
  \newcommand{\xxxSubParagraphStar}[1]{\oldsubparagraph*{#1}\mbox{}}
  \newcommand{\xxxSubParagraphNoStar}[1]{\oldsubparagraph{#1}\mbox{}}
\fi
\makeatother


\usepackage{longtable,booktabs,array}
\newcounter{none} % for unnumbered tables
\usepackage{calc} % for calculating minipage widths
% Correct order of tables after \paragraph or \subparagraph
\usepackage{etoolbox}
\makeatletter
\patchcmd\longtable{\par}{\if@noskipsec\mbox{}\fi\par}{}{}
\makeatother
% Allow footnotes in longtable head/foot
\IfFileExists{footnotehyper.sty}{\usepackage{footnotehyper}}{\usepackage{footnote}}
\makesavenoteenv{longtable}
\usepackage{graphicx}
\makeatletter
\newsavebox\pandoc@box
\newcommand*\pandocbounded[1]{% scales image to fit in text height/width
  \sbox\pandoc@box{#1}%
  \Gscale@div\@tempa{\textheight}{\dimexpr\ht\pandoc@box+\dp\pandoc@box\relax}%
  \Gscale@div\@tempb{\linewidth}{\wd\pandoc@box}%
  \ifdim\@tempb\p@<\@tempa\p@\let\@tempa\@tempb\fi% select the smaller of both
  \ifdim\@tempa\p@<\p@\scalebox{\@tempa}{\usebox\pandoc@box}%
  \else\usebox{\pandoc@box}%
  \fi%
}
% Set default figure placement to htbp
\def\fps@figure{htbp}
\makeatother


% definitions for citeproc citations
\NewDocumentCommand\citeproctext{}{}
\NewDocumentCommand\citeproc{mm}{%
  \begingroup\def\citeproctext{#2}\cite{#1}\endgroup}
\makeatletter
 % allow citations to break across lines
 \let\@cite@ofmt\@firstofone
 % avoid brackets around text for \cite:
 \def\@biblabel#1{}
 \def\@cite#1#2{{#1\if@tempswa , #2\fi}}
\makeatother
\newlength{\cslhangindent}
\setlength{\cslhangindent}{1.5em}
\newlength{\csllabelwidth}
\setlength{\csllabelwidth}{3em}
\newenvironment{CSLReferences}[2] % #1 hanging-indent, #2 entry-spacing
 {\begin{list}{}{%
  \setlength{\itemindent}{0pt}
  \setlength{\leftmargin}{0pt}
  \setlength{\parsep}{0pt}
  % turn on hanging indent if param 1 is 1
  \ifodd #1
   \setlength{\leftmargin}{\cslhangindent}
   \setlength{\itemindent}{-1\cslhangindent}
  \fi
  % set entry spacing
  \setlength{\itemsep}{#2\baselineskip}}}
 {\end{list}}
\usepackage{calc}
\newcommand{\CSLBlock}[1]{\hfill\break\parbox[t]{\linewidth}{\strut\ignorespaces#1\strut}}
\newcommand{\CSLLeftMargin}[1]{\parbox[t]{\csllabelwidth}{\strut#1\strut}}
\newcommand{\CSLRightInline}[1]{\parbox[t]{\linewidth - \csllabelwidth}{\strut#1\strut}}
\newcommand{\CSLIndent}[1]{\hspace{\cslhangindent}#1}



\setlength{\emergencystretch}{3em} % prevent overfull lines

\providecommand{\tightlist}{%
  \setlength{\itemsep}{0pt}\setlength{\parskip}{0pt}}



 


\KOMAoption{captions}{tableheading}
\makeatletter
\@ifpackageloaded{caption}{}{\usepackage{caption}}
\AtBeginDocument{%
\ifdefined\contentsname
  \renewcommand*\contentsname{Table of contents}
\else
  \newcommand\contentsname{Table of contents}
\fi
\ifdefined\listfigurename
  \renewcommand*\listfigurename{List of Figures}
\else
  \newcommand\listfigurename{List of Figures}
\fi
\ifdefined\listtablename
  \renewcommand*\listtablename{List of Tables}
\else
  \newcommand\listtablename{List of Tables}
\fi
\ifdefined\figurename
  \renewcommand*\figurename{Figure}
\else
  \newcommand\figurename{Figure}
\fi
\ifdefined\tablename
  \renewcommand*\tablename{Table}
\else
  \newcommand\tablename{Table}
\fi
}
\@ifpackageloaded{float}{}{\usepackage{float}}
\floatstyle{ruled}
\@ifundefined{c@chapter}{\newfloat{codelisting}{h}{lop}}{\newfloat{codelisting}{h}{lop}[chapter]}
\floatname{codelisting}{Listing}
\newcommand*\listoflistings{\listof{codelisting}{List of Listings}}
\makeatother
\makeatletter
\makeatother
\makeatletter
\@ifpackageloaded{caption}{}{\usepackage{caption}}
\@ifpackageloaded{subcaption}{}{\usepackage{subcaption}}
\makeatother
\usepackage{bookmark}
\IfFileExists{xurl.sty}{\usepackage{xurl}}{} % add URL line breaks if available
\urlstyle{same}
\hypersetup{
  pdftitle={Smiling your way to happiness or misery? Experimental tests of competing perspectives},
  pdfauthor={Nicholas A. Coles; Annabel V. Dang; Joao Francisco Goes Braga Takayanagi},
  pdfkeywords={emotion, emotion embodiment, emotion regulation, emotion
labor, facial feedback, happiness},
  colorlinks=true,
  linkcolor={blue},
  filecolor={Maroon},
  citecolor={Blue},
  urlcolor={Blue},
  pdfcreator={LaTeX via pandoc}}


\title{Smiling your way to happiness or misery? Experimental tests of
competing perspectives}
\author{Nicholas A. Coles \and Annabel V. Dang \and Joao Francisco Goes
Braga Takayanagi}
\date{}
\begin{document}
\maketitle
\begin{abstract}
Researchers and lay people alike have long been interested in whether
people can make themselves feel happier if they simply fake a smile.
Thus far, two largely disconnected research communities have provided
seemingly competing answers to this question, with (a) emotion
embodiment researchers observing positive causal links in experimental
contexts, and (b) emotion labor researchers observing negative
associations in real-world work contexts. In the present work, we
conduct a pre-registered conceptual extension of the largest
experimental demonstration of the positive effects of posed smiles on
mood, focusing on three factors that may explain why such effects fail
to generalize to real-world work contexts. More specifically, we test
whether posed smiles fail to improve mood (1) in negative contexts, (2)
when performed repeatedly, and/or (3) when performed to avoid
punishment.
\end{abstract}


\section{Introduction}\label{introduction}

``You'll find that life is still worthwhile, if you just smile'' -- Nat
King Cole (1954)

``Pretending to be happy is almost like being happy. Until you remember
that you're only pretending. Then you're sad. Really sad.'' -- Brittainy
C. Cherry (2014)

Researchers, artists, and lay people alike have long wondered if people
can make themselves feel happier if they simply fake a smile (Coles,
Larsen, \& Lench, 2019). In addition to its allure as a cheap and
scalable intervention for promoting happiness, such an idea connects to
fundamental questions about the nature of emotion. For over a century,
emotion embodiment researchers have posited that peripheral nervous
system activity causally shapes emotion-related processes, including the
subjective experience of emotion itself (Barrett, 2017; Cacioppo et al.,
1992; Damasio \& Carvalho, 2013; Feldman et al., 2024; James, 1884,
1894; Niedenthal, 2007; Scherer \& Moors, 2019; Tomkins, 1962; Wood et
al., 2016). To evaluate such possibilities, many researchers have
conducted experiments with an easy-to-manipulate component of the
peripheral nervous system: facial musculature. For example, in a study
ostensibly about how muscle movements influence galvanic skin response,
Coles, Gaertner, Frohlich, Larsen, and Basnight-Brown (Study 1, 2023)
found that participants reported higher levels of happiness immediately
after completing muscle movement trials that configured their face into
a smile vs.~scowl. Similar patterns have been observed in (1) multiple
pre-registered replications (Baker et al., 2025; Coles et al., 2023;
Noah et al., 2018; but see {Wagenmakers et al.}, 2016), (2) a
pre-registered meta-analysis of nearly 50 years of research (Coles,
Larsen, \& Lench, 2019), and (3) a large-scale adversarial collaboration
conducted across 19 countries ({Coles et al.}, 2022).

The emotion embodiment literature provides both a theoretical and
empirical rationale for suggesting that people can make themselves feel
happier if they ``just smile'' (Cole, 1954). However, if true, it is
remarkable that people are generally \emph{least happy} in jobs that
notoriously require such behaviors, such as food and retail services
(Smith, 2007). In actuality, researcherswho study these contexts often
describe the negative effects of forcing one's emotional expression to
match the demands of a job (Hochschild, 1979)\emph{.} Emotion labor
researchers often describe posed smiles as ``surface acting'', an act of
feigning required emotional displays without changing internal states,
which is often (imperfectly) likened to the notion of ``expressive
suppression'' from research on basic emotion regulation strategies
(e.g., Gross, 1998). Relative to strategies that involve re-interpreting
negative thoughts (i.e., cognitive reappraisal) to genuinely feel and
express required emotions, expressive suppression strategies are
generally characterized as inferior (Webb et al., 2012),
counterproductive (Fernandes \& Tone, 2021), and regulatorily costly in
terms of decreases in well-being (Dryman \& Heimberg, 2018; Hu et al.,
2014; Yu et al., 2023), social outcomes (Chervonsky \& Hunt, 2017), and
cognitive performance (Roche \& Arnold, 2018; Zhu et al., 2021).
Combined, this provides emotion labor researchers with both a
theoretical and empirical rationale for warning against ``pretending to
be happy'' (Cherry, 2014).

When two scientific perspectives compete, it is not uncommon for one
camp to fail to recognize each other's epistemic standards or legitimacy
(Galison, 1997). In the context of smiling and happiness, some
researchers may be tempted to dismiss the emotion labor perspective
altogether due to internal validity concerns. For context, emotion labor
researchers have documented a large number of negative correlations
between smiling and indicators of happiness, often in real-world work
contexts (Boemo et al., 2022; Fernandes \& Tone, 2021; Hülsheger \&
Schewe, 2011; Visted et al., 2018). However, there are many confounds to
consider. For example, negative correlations may be ``due to the
stressful situation that triggers negative emotions'' requiring people
to regulate their emotion expressions (Zapf et al., 2021, p. 165) -- not
the actual act of posing a smile. This could also potentially explain
why researchers in a different theoretical tradition, emotion
embodiment, observe positive effects of posed smiles: they
overwhelmingly rely on experimental methods that manipulate smiling
while controlling for potential confounding variables (e.g., the
emotional nature of the situation; Coles, Larsen, \& Lench (2019)).

Conversely, some researchers may be tempted to dismiss the emotion
embodiment perspective over concerns about the ecological validity of
their experimental settings. Experiments often bear little resemblance
to real-world work conditions where people pose smiles -- i.e., the
setting where such acts tend to be linked with negative emotional
outcomes. For example, consider the largest experimental demonstration
of the effects of posed smiles on happiness ({Coles et al.}, 2022).
Participants were led to believe the study was focused on multi-tasking
and were asked to complete a variety of movements (e.g., tapping their
leg). For one of those movements, participants were asked to configure
their face into a smile by following muscle-by-muscle movement
instructions. Participants did indeed report improved moods after posing
smiles. However, it's possible that such effects do not generalize
beyond these sets of relatively artificial conditions (e.g., to
real-world work contexts).

\subsection{Present Work}\label{present-work}

The present work is designed to evaluate why the effects of posed smiles
are often observed as positive in emotion embodiment research but
negative in emotion labor research. We considered multiple way that
these contexts differ\footnote{Although we limit our initial
  investigation to three differences between the typical emotion
  embodiment experiment and emotion labor context, there are several
  other important differences to later consider, including: the status
  of interacting partners and the perceived goal of posing smiles.},
ultimately choosing three to focus our initial inquiry: (1) the
emotional contexts in which participants pose smiles, (2) the number of
times they pose the smile, and (3) whether non-compliance could lead to
financial penalties.

\subsubsection{The Emotion Incongruence Hypothesis: Posed Smiles Do Not
Improve Mood in Negative
Contexts}\label{the-emotion-incongruence-hypothesis-posed-smiles-do-not-improve-mood-in-negative-contexts}

First, we consider an \emph{emotion incongruence hypothesis}. Emotion
embodiment researchers often study the effects of posed expressions in a
context that is either neutral or emotionally congruent with the posed
expression (e.g., positive if posing smiles; Coles, Larsen, \& Lench
(2019)). Once again, consider the largest experimental demonstration of
such effects ({Coles et al.}, 2022). Half of participants posed smiles
while essentially staring at a blank screen, and the other half posed
smiles while looking at mildly pleasant photos. However, when extending
similar paradigms to a stressful context -- like being submerged in ice
water -- researchers have thus far failed that posed smiles improve mood
(Kraft \& Pressman, 2012; Luu et al., 2025; Pressman et al., 2021). In
addition to providing one explanation for discrepancies between the
emotion embodiment and labor literatures, such results challenge
popularized claims about the ability for facial movement manipulations
to improve mood in everyday negative situations. This includes not only
recommendations to ``just smile'', but also the neuromuscular electrical
stimulation of facial musculature associated with positive expressions
(Efthimiou, Baker, Elsenaar, et al., 2024; Efthimiou, Baker, Clarke, et
al., 2024) or the use of Botox treatments to inhibit activation in
facial musculature associated with negative expressions (Schulze et al.,
2021; Wollmer et al., 2022). (For critiques of such clinical
translations, see Coles \& Larsen (2021); Coles, Larsen, Kuribayashi, et
al. (2019)).

\subsubsection{\texorpdfstring{\textbf{\emph{The Repetition Hypothesis:
Posed Smiles Do Not Improve Mood When Performed
Repeatedly}}}{The Repetition Hypothesis: Posed Smiles Do Not Improve Mood When Performed Repeatedly}}\label{the-repetition-hypothesis-posed-smiles-do-not-improve-mood-when-performed-repeatedly}

Second, we consider a \emph{repetition hypothesis}. According to this
hypothesis, posed smiles may improve mood when performed once or twice
-- e.g., in the context of a typical emotion embodiment experiment.
However, such effects may disappear (and even reverse) when people pose
smiles repeatedly -- e.g., in the context of a typical workplace studied
by emotion embodiment surveys. Practically, such a result would
highlight the importance of posing smiles in moderation. For example,
rather than requiring workers to feign happiness in every interaction,
worker well-being may be maximized by only asking them to do so in their
most difficult interactions.

\subsubsection{\texorpdfstring{\textbf{\emph{The Punishment Hypothesis:
Posed Smiles Do Not Improve Mood When Performed to Avoid
Punishment}}}{The Punishment Hypothesis: Posed Smiles Do Not Improve Mood When Performed to Avoid Punishment}}\label{the-punishment-hypothesis-posed-smiles-do-not-improve-mood-when-performed-to-avoid-punishment}

Third, we consider a \emph{punishment hypothesis}: posed smiles can
improve mood when performed relatively voluntarily in emotion embodiment
research -- but not when done under threat of punishment (e.g., reduced
tips, docked performance evaluations, or job loss) in emotion labor
research. More practically, such results directly connect to
conversations about eradicating emotion display requirements at work
altogether (Grandey et al., 2015). Such conversations often consider
seemingly conflicting evidence from the emotion embodiment literature:
that posed smiles can increase happiness in experimental contexts (e.g.,
Chandwani \& Sharma, 2015; Choi \& Kim, 2015). However, the punishment
account potentially explains why these effects do not generalize to work
settings: smiling isn't mood boosting when performed to avoid
punishment.

\section{Method}\label{method}

In the present work, we evaluate the emotion incongruence, repetition,
and punishment accounts side-by-side. To do so, we will extend a
paradigm used in the largest experimental demonstration of the
mood-boosting effects of posed smiles ({Coles et al.}, 2022) to include
three experimentally-manipulated factors: (1) the emotional contexts in
which participants pose smiles, (2) the number of times they pose the
smile, and (3) whether non-compliance could lead to financial penalties.

The present study features a 2 (Face Movement: pose smile, express
naturally) x 2 (Emotional Context: positive, negative) x 2 (Movement
Repetitions: one time, ten times) x 2 (Punishment Threat: present,
absent) design, with the first factor manipulated within-subjects. These
manipulations are designed to evaluate boundary conditions of the
mood-boosting effect of posed smiles -- although they cannot necessarily
explain \emph{why} those conditions matter. For example, Punishment
Threat may eliminate the mood-boosting effect of posed smiles because it
shifts the perceived goal of the smile (to avoid financial penalty
vs.~appease an experimenter). But it may also eliminate the
mood-boosting effect of posed smiles because it is intrinsically
unpleasant. Our Registered Report focuses on identifying these boundary
conditions with a pre-registered and pre-vetted design, sampling plan,
and analysis strategy
(\url{https://osf.io/2suyx/?view_only=8052e16be97d4bee8f93a6fa67b02893}).

\subsection{Participants}\label{participants}

N = 1200 participants will be recruited from Prolific: an online
platform where people complete studies in exchange for payment. (See
\emph{Power Simulation} for more information.) Workers on the platform
will be eligible to participate if they (a) speak English, (b) have a
functioning webcam, (c) have previously completed at least 5 studies,
and (d) have at least 90\% of their study submissions accepted. These
latter two criteria help ensure that the quality of participants'
responses has been vetted by other researchers.

\subsection{Design}\label{design}

\subsection{Procedure}\label{procedure}

Participants will complete the online experiment via a browser-based
platform, Gorilla Experiment Builder (Anwyl-Irvine et al., 2020). To
minimize the impact of demand characteristics (Coles et al., 2025; Orne,
1962), participants will be told the study examines ``how physical
multi-tasking influences mathematical speed and accuracy''. For payment,
participants will be told that they will receive up to \$6. However,
half of participants will be informed that \$2 will be deducted if we do
not find that they followed instructions satisfactorily during recorded
segments of the study. This serves as our manipulation of Punishment
Threat, but all participants will ultimately receive full compensation
(\$6). Regardless of assigned condition, participants will be informed
of recording and will confirm webcam functionality by submitting and
confirming the playback of a short video through the survey program.

During the study, participants will complete four trials, all of which
are conceptual replications of a large multi-lab adversarial test of
facial feedback effects ({Coles et al.}, 2022). The first and last trial
are sham tasks designed to corroborate the cover story. E.g., in the
first trial, participants will be asked to ``place your left hand behind
your head and blink your eyes once per second for 5 seconds''. In the
last trial, they will be asked to ``Tap your left leg with your
right-hand index finger once per second for 5 seconds.''

More critically, in the second and third trials, participants will be
instructed to pose smiles or express themselves normally for 5 seconds.
This serves as our manipulation of Face Movement. Participants will be
shown four images of actors trained to display prototypical expressions
of happiness. These images were originally compiled from the Extended
Cohn-Kanade Dataset (Lucey et al., 2010) and arranged in a single 2x2
image matrix ({Coles et al.}, 2022). Participants will be asked to
either (a) adopt the pictured expression, or (b) not attempt to control
their face movements. Importantly, having participants view the same
image matrix in both trials helps eliminate the potentially confounding
effect that images of smiling people may have on emotional experience.
Limiting the pose to 5 seconds helps ensure a similar duration as
authentic emotional expressions (Ekman, 2003).

During the two face movement trials, participants will encounter images
previously validated to elicit either a positive or negative mood
(Coles, 2020). This serves as our manipulation of Emotional Context. For
positive contexts, participants will view images of flowers, kittens,
and rainbows, arranged in a single 1 x 3 row. For negative contexts,
they will view images of inconsiderately parked vehicles, acts of theft,
and animal maltreatment (also arranged in a single row). Images will be
displayed alongside an on-screen timer while participants complete the
face movement task. While doing so, participants will be recorded
through their webcam, which will be automatically toggled on and off by
the survey program.

Half of participants will be randomly assigned to complete the face
movement task once before proceeding to the next trial, which serves as
a conceptual replication of Coles et al. (2022). The other half of
participants will be assigned to complete the task ten times. This
serves as our manipulation of Repetition. Immediately after completing
the face movement task (either once or ten times), participants will
complete a modified Discrete Emotions Questionnaire, wherein they will
report the extent to which they experienced a variety of emotions during
the preceding trial (1 = Not at all, 7 = An extreme amount; Harmon-Jones
et al. (2016)). Following Coles et al. (2022), we modified the
questionnaire in three ways. First, to shorten the length, we limited
ourselves to three items per subscale. Second, we replaced the original
anger items (anger, rage, mad, and pissed off) with more moderate
options (e.g., irritation). Third, we limited ourselves to three emotion
subscales: happiness (three averaged items: happiness, satisfaction, and
enjoyment), anger (irritation, aggravation, and annoyance), and fear
(alarmed, scared, fear). Multiple emotion subscales were used to obscure
our interest in any specific emotional state.

The key outcome-of-interest is self-reported happiness. However, we will
also explore effects on two other outcomes frequently discussed in the
emotion labor literature: well-being and burnout. After self-reporting
their emotions, participants will complete measures of satisfaction with
life (e.g., ``In most ways my life is close to my ideal''; 1 = Strongly
Disagree, 7 = Strongly Agree; Diener et al. (1985)) and burnout (e.g.,
``To what extent to you feel hopeless''; 1 = Not at all, 7 = An extreme
amount; Malach-Pines (2005)). Afterwards, participants will complete a
simple filler math problem (e.g., 2+ 2 = ?). This is done to both (a)
corroborate the cover story and (b) evaluate whether participants are
paying attention.

After completing all trials, participants will self-report their age,
gender (Male, Female, Nonbinary; Garbarski et al. (2025)), and ethnicity
(Revesz, 2024). Afterwards, participants will be asked two questions in
a funnel debriefing. First, participants will simply be asked what they
believe was the purpose of this study. Second, they will be informed
that ``There are details about the purpose of this study that we have
not yet told you about'' and asked to guess those details. Afterwards,
participants will be fully debriefed.

After the study is complete, we will evaluate how well participants
completed the face movement tasks. To do so, we will use a commercially
available emotion expression recognition model, Hume AI
(\url{https://www.hume.ai/}). Like other emotion expression recognition
models, Hume AI uses computer vision to track facial landmarks (e.g.,
lips) associated with emotion expression (e.g., smiling) in each frame
of a user-uploaded video. These features are passed through a
pre-trained deep neural network to predict observer ratings of emotion
(Cowen et al., 2021). To be clear, such models can fail to accurately
infer peoples' subjective emotional experiences in complex and/or
real-world scenarios (Cross et al., 2023). However, in relatively simple
experimental contexts like ours, such models have been successfully
deployed to identify non-smiling participants in an efficient and
scalable manner (Coles, 2020; {Coles et al.}, 2022). In the present
work, we will use Hume AI to generate moment-to-moment estimates of
smiling-related activity using Facial Action Coding System
classifications (Ekman \& Friesen, 1978). We will focus on Action Unit
12 in particular (measured on a 0 -- 1 scale), which is associated with
pulling the corner of the lips back towards the ear.

\subsection{Measures}\label{measures}

\section{Results}\label{results}

As a reminder, our study deploys a 2 (Face Movement: pose smile, express
naturally) x 2 (Emotional Context: positive, negative) x 2 (Movement
Repetitions: one time, ten times) x 2 (Punishment threat: present,
absent) design, with Face Movement manipulated within-subject.

The emotion incongruence, repetition, and punishment hypotheses all
predict a two-way interaction with Face Movement -- but differ in their
proposed moderating factor. More specifically, the emotion incongruence
hypothesis predicts a two-way interaction between Face Movement and
Emotional Context, wherein there are only positive simple effects of
Face Movement in positive contexts (Table 1). The repetition hypothesis
predicts a two-way interaction between Face Movement and Movement
Repetition, wherein Face Movement only has a positive simple effect when
repeated once (Table 1). Last, the punishment hypothesis predicts a
two-way interaction between Face Movement and Punishment, wherein there
are only positive simple effects of Face Movement when there is no
stated punishment for non-adherence (Table 1).

\section{Discussion}\label{discussion}

\section*{References}\label{references}
\addcontentsline{toc}{section}{References}

\protect\phantomsection\label{refs}
\begin{CSLReferences}{1}{0}
\bibitem[\citeproctext]{ref-anwylirvine2020}
Anwyl-Irvine, A. L., Massonnié, J., Flitton, A., Kirkham, N., \&
Evershed, J. K. (2020). Gorilla in our midst: An online behavioral
experiment builder. \emph{Behavior Research Methods}, \emph{52},
388--407. \url{https://doi.org/10.3758/s13428-019-01237-x}

\bibitem[\citeproctext]{ref-baker2025}
Baker, J., Van der Donck, S., Boets, B., \& Korb, S. (2025). Increasing
perceived happiness in neutral faces by posing a smile: An
electroencephalography ({EEG}) frequency-tagging study. \emph{Emotion}.
\url{https://doi.org/10.1037/emo0001514}

\bibitem[\citeproctext]{ref-barrett2017}
Barrett, L. F. (2017). The theory of constructed emotion: An active
inference account of interoception and categorization. \emph{Social
Cognitive and Affective Neuroscience}, \emph{12}(1), 1--23.
\url{https://doi.org/10.1093/scan/nsw154}

\bibitem[\citeproctext]{ref-boemo2022}
Boemo, T., Nieto, I., Vazquez, C., \& Sanchez-Lopez, A. (2022).
Relations between emotion regulation strategies and affect in daily
life: A systematic review and meta-analysis of studies using ecological
momentary assessments. \emph{Neuroscience \& Biobehavioral Reviews},
\emph{139}, 104747.
\url{https://doi.org/10.1016/j.neubiorev.2022.104747}

\bibitem[\citeproctext]{ref-cacioppo1992}
Cacioppo, J. T., Berntson, G. G., \& Klein, D. J. (1992). \emph{Emotion
and social behavior}. Sage Publications, Inc.

\bibitem[\citeproctext]{ref-chandwani2015}
Chandwani, R., \& Sharma, D. (2015). \emph{Managing emotions: Emotional
labor or emotional enrichment}. Indian Institute of Management.

\bibitem[\citeproctext]{ref-cherry2014}
Cherry, C. B. (2014). \emph{Loving {Mr. Daniels}}. CreateSpace
Independent Publishing Platform.

\bibitem[\citeproctext]{ref-chervonsky2017}
Chervonsky, E., \& Hunt, C. (2017). Suppression and expression of
emotion in social and interpersonal outcomes: A meta-analysis.
\emph{Emotion}, \emph{17}(4), 669--683.
\url{https://doi.org/10.1037/emo0000270}

\bibitem[\citeproctext]{ref-choi2015}
Choi, Y. G., \& Kim, K. S. (2015). A literature review of emotional
labor and emotional labor strategies. \emph{Universal Journal of
Management}, \emph{3}(7), 283--290.
\url{https://doi.org/10.13189/ujm.2015.030704}

\bibitem[\citeproctext]{ref-cole1954}
Cole, N. K. (1954). \emph{Smile}. Song.

\bibitem[\citeproctext]{ref-coles2020}
Coles, N. A. (2020). \emph{Face and feeling: An examination of the role
of facial feedback in emotion} {[}PhD thesis, University of
Tennessee{]}. \url{https://trace.tennessee.edu/utk_graddiss/5874}

\bibitem[\citeproctext]{ref-coles2023artifact}
Coles, N. A., Gaertner, L., Frohlich, B., Larsen, J. T., \&
Basnight-Brown, D. M. (2023). Fact or artifact? Demand characteristics
and participants' beliefs can moderate, but do not fully account for,
the effects of facial feedback on emotional experience. \emph{Journal of
Personality and Social Psychology}, \emph{124}(2), 287.
\url{https://doi.org/10.1037/pspa0000316}

\bibitem[\citeproctext]{ref-coles2021eyebrows}
Coles, N. A., \& Larsen, J. T. (2021). Claims about the effects of
botulinum toxin on depression should raise some eyebrows. \emph{Journal
of Psychiatric Research}, \emph{140}, 551--552.
\url{https://doi.org/10.1016/j.jpsychires.2021.05.021}

\bibitem[\citeproctext]{ref-coles2019botox}
Coles, N. A., Larsen, J. T., Kuribayashi, J., \& Kuelz, A. (2019). Does
blocking facial feedback via botulinum toxin injections decrease
depression? A critical review and meta-analysis. \emph{Emotion Review},
\emph{11}(4), 294--309. \url{https://doi.org/10.1177/1754073919868762}

\bibitem[\citeproctext]{ref-coles2019meta}
Coles, N. A., Larsen, J. T., \& Lench, H. C. (2019). A meta-analysis of
the facial feedback literature: Effects of facial feedback on emotional
experience are small and variable. \emph{Psychological Bulletin},
\emph{145}(6), 610. \url{https://doi.org/10.1037/bul0000194}

\bibitem[\citeproctext]{ref-coles2022}
{Coles, N. A., March, D. S., Marmolejo-Ramos, F., Larsen, J. T., Arinze,
N. C., Ndukaihe, I. L., et al.} (2022). A multi-lab test of the facial
feedback hypothesis by the {Many Smiles Collaboration}. \emph{Nature
Human Behaviour}, \emph{6}(12), 1731--1742.
\url{https://doi.org/10.1038/s41562-022-01458-9}

\bibitem[\citeproctext]{ref-coles2025}
Coles, N. A., Wyatt, M., \& Frank, M. C. (2025). A meta-analysis of the
impact and heterogeneity of explicit demand characteristics.
\emph{Collabra: Psychology}, \emph{11}(1), 143005.
\url{https://doi.org/10.1525/collabra.143005}

\bibitem[\citeproctext]{ref-cowen2021}
Cowen, A. S., Keltner, D., Schroff, F., Jou, B., Adam, H., \& Prasad, G.
(2021). Sixteen facial expressions occur in similar contexts worldwide.
\emph{Nature}, \emph{589}(7841), 251--257.
\url{https://doi.org/10.1038/s41586-020-3037-7}

\bibitem[\citeproctext]{ref-cross2023}
Cross, M. P., Acevedo, A. M., \& Hunter, J. F. (2023). A critique of
automated approaches to code facial expressions: What do researchers
need to know? \emph{Affective Science}, \emph{4}(3), 500--505.
\url{https://doi.org/10.1007/s42761-023-00195-0}

\bibitem[\citeproctext]{ref-damasio2013}
Damasio, A., \& Carvalho, G. B. (2013). The nature of feelings:
Evolutionary and neurobiological origins. \emph{Nature Reviews
Neuroscience}, \emph{14}(2), 143--152.
\url{https://doi.org/10.1038/nrn3403}

\bibitem[\citeproctext]{ref-diener1985}
Diener, E., Emmons, R. A., Larsen, R. J., \& Griffin, S. (1985). The
satisfaction with life scale. \emph{Journal of Personality Assessment},
\emph{49}(1), 71--75. \url{https://doi.org/10.1207/s15327752jpa4901_13}

\bibitem[\citeproctext]{ref-dryman2018}
Dryman, M. T., \& Heimberg, R. G. (2018). Emotion regulation in social
anxiety and depression: A systematic review of expressive suppression
and cognitive reappraisal. \emph{Clinical Psychology Review}, \emph{65},
17--42. \url{https://doi.org/10.1016/j.cpr.2018.07.004}

\bibitem[\citeproctext]{ref-efthimiou2024happiness}
Efthimiou, T.-N., Baker, J., Clarke, A., Elsenaar, A., Mehu, M., \&
Korb, S. (2024). Zygomaticus activation through facial neuromuscular
electrical stimulation ({fNMES}) induces happiness perception in
ambiguous facial expressions and affects neural correlates of face
processing. \emph{Social Cognitive and Affective Neuroscience},
\emph{19}(1), nsae013. \url{https://doi.org/10.1093/scan/nsae013}

\bibitem[\citeproctext]{ref-efthimiou2024smiling}
Efthimiou, T.-N., Baker, J., Elsenaar, A., Mehu, M., \& Korb, S. (2024).
Smiling and frowning induced by facial neuromuscular electrical
stimulation ({fNMES}) modulate felt emotion and physiology.
\emph{Emotion}. \url{https://doi.org/10.1037/emo0001408}

\bibitem[\citeproctext]{ref-ekman2003}
Ekman, P. (2003). Darwin, deception, and facial expression. \emph{Annals
of the New York Academy of Sciences}, \emph{1000}(1), 205--221.
\url{https://doi.org/10.1196/annals.1280.010}

\bibitem[\citeproctext]{ref-ekman1978}
Ekman, P., \& Friesen, W. V. (1978). \emph{Facial action coding system}.
Environmental Psychology \& Nonverbal Behavior.
\url{https://doi.org/10.1037/t27734-000}

\bibitem[\citeproctext]{ref-feldman2024}
Feldman, M., Bliss-Moreau, E., \& Lindquist, K. (2024). The neurobiology
of interoception and affect. \emph{Trends in Cognitive Sciences}.
\url{https://doi.org/10.1016/j.tics.2024.01.009}

\bibitem[\citeproctext]{ref-fernandes2021}
Fernandes, M. A., \& Tone, E. B. (2021). A systematic review and
meta-analysis of the association between expressive suppression and
positive affect. \emph{Clinical Psychology Review}, \emph{88}, 102068.
\url{https://doi.org/10.1016/j.cpr.2021.102068}

\bibitem[\citeproctext]{ref-galison1997}
Galison, P. (1997). \emph{Image and logic: A material culture of
microphysics}. University of Chicago Press.

\bibitem[\citeproctext]{ref-garbarski2025}
Garbarski, D., Dykema, J., Yonker, J. A., Bae, R. E., \& Rosenfeld, R.
A. (2025). Improving the measurement of gender in surveys: Effects of
categorical versus open-ended response formats on measurement and data
quality among college students. \emph{Journal of Survey Statistics and
Methodology}, \emph{13}(1), 18--38.
\url{https://doi.org/10.1093/jssam/smae043}

\bibitem[\citeproctext]{ref-grandey2015}
Grandey, A. A., Rupp, D., \& Brice, W. N. (2015). Emotional labor
threatens decent work: A proposal to eradicate emotional display rules.
\emph{Journal of Organizational Behavior}, \emph{36}(6), 770--785.
\url{https://doi.org/10.1002/job.2020}

\bibitem[\citeproctext]{ref-gross1998}
Gross, J. J. (1998). The emerging field of emotion regulation: An
integrative review. \emph{Review of General Psychology}, \emph{2}(3),
271--299. \url{https://doi.org/10.1037/1089-2680.2.3.271}

\bibitem[\citeproctext]{ref-harmonjones2016}
Harmon-Jones, C., Bastian, B., \& Harmon-Jones, E. (2016). The discrete
emotions questionnaire: A new tool for measuring state self-reported
emotions. \emph{PloS One}, \emph{11}(8), e0159915.
\url{https://doi.org/10.1371/journal.pone.0159915}

\bibitem[\citeproctext]{ref-hochschild1979}
Hochschild, A. R. (1979). Emotion work, feeling rules, and social
structure. \emph{American Journal of Sociology}, \emph{85}(3), 551--575.
\url{https://www.jstor.org/stable/2778583}

\bibitem[\citeproctext]{ref-hu2014}
Hu, T., Zhang, D., Wang, J., Mistry, R., Ran, G., \& Wang, X. (2014).
Relation between emotion regulation and mental health: A meta-analysis
review. \emph{Psychological Reports}, \emph{114}(2), 341--362.
\url{https://doi.org/10.1037/a0027600}

\bibitem[\citeproctext]{ref-hulsheger2011}
Hülsheger, U. R., \& Schewe, A. F. (2011). On the costs and benefits of
emotional labor: A meta-analysis of three decades of research.
\emph{Journal of Occupational Health Psychology}, \emph{16}(3), 361.
\url{https://doi.org/10.1037/a0022876}

\bibitem[\citeproctext]{ref-james1884}
James, W. (1884). What is emotion? \emph{Mind}, \emph{8}, 188--205.
\url{https://www.jstor.org/stable/2246769}

\bibitem[\citeproctext]{ref-james1894}
James, W. (1894). Discussion: The physical basis of emotion.
\emph{Psychological Review}, \emph{1}, 516--529.
\url{https://doi.org/10.1037/h0065078}

\bibitem[\citeproctext]{ref-kraft2012}
Kraft, T. L., \& Pressman, S. D. (2012). Grin and bear it: The influence
of manipulated facial expression on the stress response.
\emph{Psychological Science}, \emph{23}(11), 1372--1378.
\url{https://doi.org/10.1177/0956797612445312}

\bibitem[\citeproctext]{ref-lucey2010}
Lucey, P., Cohn, J. F., Kanade, T., Saragih, J., Ambadar, Z., \&
Matthews, I. (2010). The extended {Cohn-Kanade} dataset ({CK+}): A
complete dataset for action unit and emotion-specified expression.
\emph{2010 IEEE Computer Society Conference on Computer Vision and
Pattern Recognition - Workshops}, 94--101.
\url{https://doi.org/10.1109/CVPRW.2010.5543262}

\bibitem[\citeproctext]{ref-luu2025}
Luu, J. H., Acevedo, A. M., Pourmand, V., \& Pressman, S. D. (2025). The
power of smiles: Mitigating pain through facial expression. \emph{The
Journal of Positive Psychology}, 1--10.
\url{https://doi.org/10.1080/17439760.2025.2462231}

\bibitem[\citeproctext]{ref-malachpines2005}
Malach-Pines, A. (2005). The burnout measure, short version.
\emph{International Journal of Stress Management}, \emph{12}(1), 78--88.
\url{https://doi.org/10.1037/1072-5245.12.1.78}

\bibitem[\citeproctext]{ref-niedenthal2007}
Niedenthal, P. M. (2007). Embodying emotion. \emph{Science},
\emph{316}(5827), 1002--1005.
\url{https://doi.org/10.1126/science.1136930}

\bibitem[\citeproctext]{ref-noah2018}
Noah, T., Schul, Y., \& Mayo, R. (2018). When both the original study
and its failed replication are correct: Feeling observed eliminates the
facial-feedback effect. \emph{Journal of Personality and Social
Psychology}, \emph{114}(5), 657.
\url{https://doi.org/10.1037/pspa0000121}

\bibitem[\citeproctext]{ref-orne1962}
Orne, M. T. (1962). On the social psychology of the psychological
experiment: With particular reference to demand characteristics and
their implications. \emph{American Psychologist}, \emph{17}(11),
776--783. \url{https://doi.org/10.1037/h0043424}

\bibitem[\citeproctext]{ref-pressman2021}
Pressman, S. D., Acevedo, A. M., Hammond, K. V., \& Kraft-Feil, T. L.
(2021). Smile (or grimace) through the pain? The effects of
experimentally manipulated facial expressions on needle-injection
responses. \emph{Emotion}, \emph{21}(6), 1188.
\url{https://doi.org/10.1037/emo0000913}

\bibitem[\citeproctext]{ref-revesz2024}
Revesz, R. L. (2024). \emph{Revisions to {OMB}'s statistical policy
directive no. 15: Standards for maintaining, collecting, and presenting
federal data on race and ethnicity}. Federal Register.

\bibitem[\citeproctext]{ref-roche2018}
Roche, J. M., \& Arnold, H. S. (2018). The effects of emotion
suppression during language planning and production. \emph{Journal of
Speech, Language, and Hearing Research}, \emph{61}(8), 2076--2083.
\url{https://doi.org/10.1044/2018_JSLHR-L-17-0232}

\bibitem[\citeproctext]{ref-scherer2019}
Scherer, K. R., \& Moors, A. (2019). The emotion process: Event
appraisal and component differentiation. \emph{Annual Review of
Psychology}, \emph{70}(1), 719--745.
\url{https://doi.org/10.1146/annurev-psych-122216-011854}

\bibitem[\citeproctext]{ref-schulze2021}
Schulze, J., Neumann, I., Magid, M., Finzi, E., Sinke, C., Wollmer, M.
A., \& Krüger, T. H. (2021). Botulinum toxin for the management of
depression: An updated review of the evidence and meta-analysis.
\emph{Journal of Psychiatric Research}, \emph{135}, 332--340.
\url{https://doi.org/10.15698/mic2021.02.741}

\bibitem[\citeproctext]{ref-smith2007}
Smith, T. W. (2007). \emph{Job satisfaction in the {United States}}.
National Opinion Research Center, University of Chicago.

\bibitem[\citeproctext]{ref-tomkins1962}
Tomkins, S. (1962). \emph{Affect imagery consciousness: Volume {I}: The
positive affects}. Springer Publishing Company.

\bibitem[\citeproctext]{ref-visted2018}
Visted, E., Vøllestad, J., Nielsen, M. B., \& Schanche, E. (2018).
Emotion regulation in current and remitted depression: A systematic
review and meta-analysis. \emph{Frontiers in Psychology}, \emph{9}, 756.
\url{https://doi.org/10.3389/fpsyg.2018.00756}

\bibitem[\citeproctext]{ref-wagenmakers2016}
{Wagenmakers, E.-J., Beek, T., Dijkhoff, L., Gronau, Q. F., Acosta, A.,
Adams Jr, R. B., et al.} (2016). Registered replication report: {Strack,
Martin, \& Stepper (1988)}. \emph{Perspectives on Psychological
Science}, \emph{11}(6), 917--928.
\url{https://doi.org/10.1177/1745691616674458}

\bibitem[\citeproctext]{ref-webb2012}
Webb, T. L., Miles, E., \& Sheeran, P. (2012). Dealing with feeling: A
meta-analysis of the effectiveness of strategies derived from the
process model of emotion regulation. \emph{Psychological Bulletin},
\emph{138}(4), 775. \url{https://doi.org/10.1037/a0027600}

\bibitem[\citeproctext]{ref-wollmer2022}
Wollmer, M. A., Magid, M., Kruger, T. H., \& Finzi, E. (2022). Treatment
of depression with botulinum toxin. \emph{Toxins}, \emph{14}(6), 383.
\url{https://doi.org/10.3390/toxins14060383}

\bibitem[\citeproctext]{ref-wood2016}
Wood, A., Rychlowska, M., Korb, S., \& Niedenthal, P. (2016). Fashioning
the face: Sensorimotor simulation contributes to facial expression
recognition. \emph{Trends in Cognitive Sciences}, \emph{20}(3),
227--240. \url{https://doi.org/10.1016/j.tics.2015.12.010}

\bibitem[\citeproctext]{ref-yu2023}
Yu, C.-W. F., Haase, C. M., \& Chang, J.-H. (2023). Habitual expressive
suppression of positive, but not negative, emotions consistently
predicts lower well-being across two culturally distinct regions.
\emph{Affective Science}, \emph{4}(4), 684--701.
\url{https://doi.org/10.1007/s42761-023-00221-1}

\bibitem[\citeproctext]{ref-zapf2021}
Zapf, D., Kern, M., Tschan, F., Holman, D., \& Semmer, N. K. (2021).
Emotion work: A work psychology perspective. \emph{Annual Review of
Organizational Psychology and Organizational Behavior}, \emph{8}(1),
139--172. \url{https://doi.org/10.1146/annurev-orgpsych-012420-062451}

\bibitem[\citeproctext]{ref-zhu2021}
Zhu, C., Jiang, Y., Wang, Y., Liu, D., \& Luo, W. (2021). Arithmetic
performance is modulated by cognitive reappraisal and expression
suppression: Evidence from behavioral and {ERP} findings.
\emph{Neuropsychologia}, \emph{162}, 108060.
\url{https://doi.org/10.1016/j.neuropsychologia.2021.108060}

\end{CSLReferences}




\end{document}
